\subsection*{The mathematical core}
Use $n \ge 1$ unit squares, without overlap, connected by full shared sides. Count all exposed unit edges, including hole edges. If $e$ is the number of shared sides, then $\mathbf{P=4n-2e}$: each meeting removes two edge contributions. Hence every perimeter is even. Adding a tile touching $k$ old sides changes $P$ by $4-2k$.

\textbf{Maximum:} Connectivity gives $e \ge n-1$, so $P \le 2n+2$; a straight row attains it. Equality holds exactly when the tile-adjacency graph is a tree. This rules out a 2-by-2 block. Whether it rules out a hole depends on corners. The Grades 4--5 page says two tiles touching at a corner leave no gap; then a tree can enclose a hole without adding a shared side. If instead a corner gap let the inside out, a hole would need a ring of side-joined tiles, a cycle, so no tree could have one (page 14).

\textbf{Minimum:} If $r$ rows and $c$ columns are occupied, then $P \ge 2(r+c)$ and $n \le rc$. For $P=2s$, the largest possible area is floor$(s/2)\cdot$ceil$(s/2)$, attained by a balanced rectangle. For fixed $n$, the minimum perimeter is 2 ceil$(2\sqrt{n})$. Near-square rectangles with corner-started partial strips attain the bound for every $n$. Equality in $P \ge 2(r+c)$ requires one occupied interval in every occupied row and column, not necessarily a full rectangle.

\subsection*{What the activity develops}
\textbf{K--1:} Build legal shapes, count exposed sides and compare records. The shared launch models both a whole-side join and boundary counting with three tiles.

\textbf{Grades 2--3:} Track shared-side cancellation and one-tile changes, then explain a longest-boundary construction.

\textbf{Grades 4--5:} Pair row/column lower bounds with attaining shapes and generalize exact area--perimeter limits. Square-root notation is an optional adult shorthand.

\subsection*{From trials to explanation}
A best shape found so far is a record and a candidate optimum. A conjecture predicts it cannot be beaten. A matching lower or upper bound, together with a construction attaining it, proves the optimum. The detailed counting arguments remain on printed page 3.
